\documentclass[10pt,a4paper]{amsart}
\usepackage[margin=19mm]{geometry}
\usepackage{amsmath,amssymb,mathtools}
\usepackage[scr=rsfs,cal=euler]{mathalfa}
\usepackage{xfrac}
\usepackage[unicode,hidelinks,bookmarks=false]{hyperref}
\newcommand{\ie}{i.e.\ }
\newcommand{\cf}{cf.\ }
\newcommand{\resp}{respectively\ }
\newcommand{\Subsection}{\subsection}
\providecommand{\nobreakdash}{\nobreak\hbox{-}}
\setlength{\emergencystretch}{2em}
\multlinegap0pt
\usepackage{enumerate}
\usepackage{amssymb}
\usepackage{cancel}
\usepackage{mathtools}
\newtheorem{theorem}{Theorem}
\newcommand{\PG}{\mathrm{PG}}
 \newcommand{\set}[1]{ \left \{ #1 \right \} }
\title{Cones from maximum $h$-scattered linear sets and a stability result for cylinders from hyperovals}
\author{Sam Adriaensen, Jonathan Mannaert, Paolo Santonastaso, Ferdinando Zullo}
\date{Source: arXiv:2210.09645v1 (CC BY 4.0)}
\begin{document}
\maketitle

\noindent\emph{Source and licence.} Cones from maximum $h$-scattered linear sets and a stability result for cylinders from hyperovals. Version arXiv:2210.09645v1 (CC BY 4.0), distributed by arXiv under the Creative Commons Attribution 4.0 International licence, http://creativecommons.org/licenses/by/4.0/, as recorded in the arXiv licence metadata for this exact version and shown on the abstract page https://arxiv.org/abs/2210.09645v1. Full licence notice: \texttt{licenses/arxiv-2210.09645-CC-BY-4.0.txt}.\par\smallskip

\noindent\textit{Editorial scope.} Counted unit: the article's theorem th:intersectionspaces (source0.tex lines 1165--1172). The article's complete original proof of it is delivered with it (source0.tex lines 1173--1209, one \texttt{proof} environment of 37 lines). No mathematics is written, repaired or re-derived: the statement and the proof are the article's own lines, in the article's own notation, and no retained line is substituted or re-flowed.  Retained uncounted as the source-local context the proof needs: lines 1080--1092.  Nothing was omitted from any retained span, so the package carries no omission record.  No retained line contains a citation, so the wrapper carries no bibliography and no bibliography entry.  No retained line contains a figure environment, an image-inclusion command or drawing code, so no image is delivered and the package declares no asset dependency.  Editorial formatting only, in addition: an AMS \texttt{amsart} preamble that restates the article's own declarations and macros used by the retained lines, namely the environments \texttt{theorem} on separate counters, together with the standard packages those lines need.  The retained environments print in this document as Theorem 1 in order of appearance, whereas the article prints the same items with the numbering of its own full document.  The complete native source of the article as distributed in the arXiv e-print for this exact version is bundled unchanged as \texttt{sources/132217/source0.tex}.
\begin{theorem} \label{th:propgenkmarcspace}
Let $S \subseteq \PG(r,q)$, with $r\geq 4$ and $q>2$, be a set of points of type $\{0,q^{r-2}+2q^{r-3},q^{r-2}+t\}_{r-1}$  of size $q^{r-1}+q^{r-2}+t$, with $2q^{r-3}< t \leq q^{r-2}+q-1$. Then the following properties are true:
\begin{itemize}
    \item[(i)] if $\pi$ is a $k$-space with $2 \leq k < r$, then $\pi \cap S = \emptyset$, or $|\pi \cap S| \geq q^{k-1}+2q^{k-2}$;
    \item[(ii)] there are no tangent lines to $S$;
    \item[(iii)] there exists at least one $k$-space that is $(q^{k-1}+2q^{k-2})$-secant, for every $k \in \{2,\ldots,r-1\}$, and at least one $2$-secant line; 
    \item[(iv)] any $(q+2)$-secant plane to $S$ meets $S$ in a hyperoval;
    \item[(v)] it holds that $q$ is even;
    \item[(vi)] all hyperplanes through an $(q^{r-3}+2q^{r-4})$-secant $(r-2)$-space are $(q^{r-2}+2q^{r-3})$-secant spaces; 
    \item[(vii)] it holds that $t=q^{r-2}$;
    \item[(viii)] every $(r-i+1)$-space through a fixed $(q^{r-i-1}+2q^{r-i-2})$-secant $(r-i)$-space is a $(q^{r-i}+2q^{r-i-1})$-secant space, for any $i \in \{1,\ldots,r-2\}$. 
\end{itemize}
\end{theorem}

\begin{theorem} 
\label{th:intersectionspaces}
Suppose that $q\geq 4$ and let $S \subseteq \PG(r,q)$ be a set of points of type $\{0,q^{r-2}+2q^{r-3},2q^{r-2}\}_{r-1}$ of size $q^{r-1}+2q^{r-2}$.
If $\tau$ is an $(r-i)$-space, with $1 \leq i \leq r-2$, then
\[
 |\tau \cap S| \in \{2q^{r-i-1}+c(-q^{r-i-1}+2q^{r-i-2}) \colon c \in \mathbb{Z}, -q \leq c \leq 1\} \cup \set 0.
\]
\end{theorem}

\begin{proof}
Let's start by proving the result when $i=2$. Let $\tau$ be a $j$-secant $(r-2)$-space. Let $m$ be the number of hyperplanes through $\tau$ of $\PG(r,q)$ that are $(q^{r-2}+2q^{r-3})$-secant. So, by counting arguments, we find that
\[
m(q^{r-2}+2q^{r-3}-j)+(q+1-m)(2q^{r-2}-j)+j=q^{r-1}+2q^{r-2},
\]
from which it follows that
\[
m(-q^{r-3}+2q^{r-4})=-q^{r-2}+j.
\]
This means that $-q^{r-3}+2q^{r-4}$ divides $-q^{r-2}+j$, thus $-q^{r-3}+2q^{r-4}$ divides $-2q^{r-3}+j$. Therefore, there exists an integer $c$ such that $j=2q^{r-3}+c(-q^{r-3}+2q^{r-4})$. By (i) of Theorem \ref{th:propgenkmarcspace} we also have 
that $c \leq 1$.\\
Now we use an inductive argument.
Suppose that the assertion is true for all the $(r-g)$-spaces with $g\leq i$, for some $i \in \{2,\ldots,r-3\}$. 
Then we want to prove the assertion for $g=i+1$.
In particular, for every secant $(r-i+1)$-spaces, its intersection number with $S$ is in the set
 $\{2q^{r-i}+c(-q^{r-i}+2q^{r-i-1}) \colon c \in \mathbb{Z}, c \leq 1\}$, and all the secant $(r-i)$-space have possible intersection numbers in the set $\{2q^{r-i-1}+c(-q^{r-i-1}+2q^{r-i-2}) \colon c \in \mathbb{Z}, c \leq 1\} $.
We claim that 
\[
 |\tau \cap S| \in \{2q^{r-i-2}+c(-q^{r-i-2}+2q^{r-i-3}) \colon c \in \mathbb{Z}, c \leq 1\}, 
\]
for every secant $(r-i-1)$-space $\tau$. First, assume that $\tau$ is a $j$-secant $(r-i-1)$-space. This means that there exists some $\overline{c} \in \mathbb{Z}$ such that $\tau$ is contained in a $(2q^{r-i}+\overline{c}(-q^{r-i}+2q^{r-i-1}))$-secant $(r-i+1)$-space $\sigma$. 
Denote the number of $(r-i)$-spaces through $\tau$ that are $(2q^{r-i-1}+\alpha(-q^{r-i-1}+2q^{r-i-2}))$-secant and contained in $\sigma$ by $m_\alpha$. As before, by counting the size of the intersection between $(r-i)$-spaces through $\tau$ and $S$ we obtain
 \[
 \sum_{\alpha} m_{\alpha} (2q^{r-i-1}+{\alpha}(-q^{r-i-1}+2q^{r-i-2})-j)+j=2q^{r-i}+\overline{c}(-q^{r-i}+2q^{r-i-1}), 
 \]
and since $\sum_{\alpha} m_{\alpha}=q+1$, we have that
\[
\begin{array}{rl}
qj & = 2q^{r-i-1}+\sum_{\alpha} m_{\alpha} {\alpha}(-q^{r-i-1}+2q^{r-i-2})-\overline{c}(-q^{r-i}+2q^{r-i-1}) \\
& = 2q^{r-i-1}+(\sum_{\alpha} m_{\alpha} {\alpha}-\overline{c}q)(-q^{r-i-1}+2q^{r-i-2})
\end{array}\] 
and so
\[
j=2q^{r-i-2}+a(-q^{r-i-2}+2q^{r-i-3}),
\]
for some integer $-q\leq a \leq 1$. 
\end{proof}

\end{document}
